\section[Les attentes fonctionnelles]{Les attentes fonctionnelles : recréer numériquement le confort du livre papier}
Dans une époque où l'urgence des moteurs imposent de produire plus en détrimant de la qualité, le confort de lecture ne doit pas être perçu comme un luxe cosmétique, mais comme un socle de la \enquote{Slow Science}. Les attentes des usagers sont donc de pouvoir recréer le confort papier du livre en le développant numérique à travers les tables des matières cliquables et le développement de l'OCR. Actuellement, l'un des avantages de CujasNum pour les usagers qui en sont \enquote{grandement satisfait est la présence de la table des matières cliquables, implémentées dans leur visionneuse JPEG mais aussi dans le PDF}. La mise en place de ce service permet aux utilisateurs de mieux appréhender le défilement fastidieux des corpus volumineux. La table des matières permet comme l'expliquait un des enseignants-chercheurs d'aller directement à la section souhaitée, complété par l'OCR. Cela favorise la recherche numérique. L'OCR a été massivement adopté par les institutions patrimoniales permettant aux chercheurs, aux étudiants de faire de la recherche plein-texte sur les images numérisées. C'est un outil très précieux et utile encrer dans les pratiques de recherche d'aujourd'hui. L'absence d'océrisation condamne le portail documentaire ou le mauvais référencement de celui-ci empêche celui-ci de pouvoir être exploité comme il le devrait. 

Les chercheurs demandent s'ils seraient possible de pouvoir mettre en place pour la future bibliothèque numérique de l'OCR qui soit capable aussi bien de faire des recherches plein-texte que de traiter des occurences précises en latin également avec le surlignage de termes lorsqu'ils sont trouvés dans la liste de résultat, accompagnée d'un lien cliquable pour se rendre directement au passage identifié, ou dans la visionneuse. L'adoption massive du standard IIIF est devenu le levier de la transformation numérique. La visionnneuse doit permettre de réaliser \enquote{un zoom de haute résolution, la comparaison des documents et des outils de captures dynamiques}. Par exemple, l'ajout des outils de captures d'écran sur la visionneuse permettrait à l'un des doctorants interviewés de mieux préparer ces illustrations pédagogiques qu'il utilise sur ces powerpoints. L'ajout de visionneuse IIIF, comme \textit{Mirador} ou \textit{Universal Viewer} a été adopté par de nombreuses institutions patrimoniales au cours de ces dernières décennies permettant à ceux qui l'utilise de pouvoir récupérer le manifeste IIIF de le consulter sur Mirador en le comparant à une autre édition présent dans une autre bibliothèque. L'avantage est que la personne n'a pas besoin de se déplacer pour cela dans l'institution. 

La préservation de la sérendipité est importante pour les usagers. Inspiré par la bibliothèque de Munich, la nouvelle bibliothèque numérique de Cujas doit proposer une interface relationnelle. La \enquote{recherche bucolique} pronée par l'un des enseignants-chercheurs n'est pas une errance, mais une méthodologie historique. Le système doit suggérer des \enquote{chemins de traverses} à travers des liens vers des éditions similaires, des auteurs cités ou des documents connexes, permettant une exploration par cercles concentriques plutôt qu'une simple réponse à une requête unique. La contrainte de la recherche n'est pas quelque chose d'appréciée par les chercheurs, c'est le risque lorsqu'il fait une requête et que celle-ci retourne trop de résultats ou alors très peu. Si la vision des collections des experts métiers s'impose de trop dans la méthode de recherche concernant la vision des collections qu'elle a pour les utilisateurs \enquote{les empêchant ainsi de faire leur propre découverte}. Cela ne lui laisse alors ni l'opportunité de s'approprier les collections par plusieurs tentatives de requêtes ni de faire des découvertes inattendues. Il est important de trouver un équilibre entre \enquote{la recherche ciblée par un but précis (trouvabilité) à la recherche exploratoire par sérendipité (découvrabilité)}, comme l'a expliqué Elina Leblanc dans sa thèse. 

Un autre garant de la sérendipité numérique est le \enquote{browsing}(le feuilletage libre), recréant l'expérience physique de la déambulation dans les rayons de la bibliothèque Cujas. Face à l'augmentation des collections numérisées, dans le monde des bibliothèques numériques, constitue \enquote{à la fois un avantage pour les lecteurs, leur offrant accès à toujours plus de documents et de connaissances. Le risque est la difficulté de trouver efficacement les ressources recherchées}. L'une des solutions est le recours au \enquote{browsing} permettant à l'utilisateur de naviguer \enquote{au milieu des collections, en empruntant des chemins qui ne lui ont pas été indiqués par la bibliothèque}. L'utilisateur peut également faire un tri dans les collections en utilisant les facettes qui lui permettent de naviguer parmi les contenus. \enquote{En choisissant lui-même les facettes, l'utilisateur crée un parcours de recherche personnalisé, qui lui permet d'investir complètement l'interface et l'encourage à poursuivre sa navigation}. L'amélioration de la recherche et du confort de lecture entraînent la nécessité de nouveaux outils permettant aux utilisateurs de sélections ce qui leur convient s'appropriant ainsi mieux le matériel. 